% !TEX root = ../main.tex

%----------------------------------------------------------------------------------------
% ABSTRACT PAGE
%----------------------------------------------------------------------------------------
\begin{abstract}
\addchaptertocentry{\abstractname}

\textbf{Resum:} \hfill

\vspace{0.3cm}

Aquest projecte presenta una plataforma web de visualització de dades hídriques de Catalunya que consolida informació de 9 embassaments i més de 200 estacions meteorològiques provinent de les APIs públiques de la Generalitat de Catalunya (Socrata). El sistema combina un pipeline ETL automatitzat amb Node.js que s'executa diàriament mitjançant GitHub Actions amb un frontend interactiu basat en React que ofereix quatre panells de visualització: mapa geogràfic amb Leaflet, evolució temporal amb Recharts, correlació precipitació-nivells, i matriu d'alertes de sequera amb indicadors semàfor i tendències. L'arquitectura desacoblada ETL/frontend permet un funcionament de cost zero aprofitant GitHub Pages i GitHub Actions. El projecte és completament de codi obert i documentat.

\vspace{1cm}

\textbf{Abstract:}

\vspace{0.3cm}

This project presents a web platform for visualizing water resource data from Catalonia, consolidating information from 9 reservoirs and over 200 meteorological stations sourced from the Catalan Government's public APIs (Socrata). The system combines an automated ETL pipeline built with Node.js that runs daily via GitHub Actions with an interactive React-based frontend featuring four visualization panels: a geographic map with Leaflet, temporal evolution with Recharts, precipitation-level correlation, and a drought alert matrix with traffic light indicators and trends. The decoupled ETL/frontend architecture enables zero-cost operation leveraging GitHub Pages and GitHub Actions. The project is fully open-source and documented.

\end{abstract}
